\section{Limits and discussion}\label{sec:limits}

\paragraph{What the language does not prevent.}  QAM secures form, not truth.  It did not prevent the dependence on a false external proof; an
audit round found it.  It does not check whether a literature value is read correctly, whether a premise holds, or whether a physical reading is
sensible.  Its schemata admit trivial models, so a record can be well-typed and empty.  And a well-typed record can still be wrong in every
statement that is not about bookkeeping.

\paragraph{Cost.}  The discipline is expensive.  The research text is dominated by protocols, and a reader who wants the results needs a separate,
short account.  Status tags and witness typing slow down writing.  In the case study the cost was worth paying, because the alternative was a long
record that nobody, including the author, could trust; whether it is worth paying in other settings is open.

\paragraph{What helped most.}  In order of effect: pre-registration with time stamps, append-only rounds, recomputation with independent code,
the review of manuscripts by a model instance without access to the record, and the typing rules.  The typing rules were most useful for the class
of error they were designed for, frame mixing, and least useful for errors of premise.

\paragraph{Next steps.}  Progress for \textsc{store}, \textsc{link} and \textsc{rollback} under explicit guard conditions; the clone and merge
instructions; a converter that adds \texttt{qam.sty} declarations to an existing research text, so that the compiler checks it; and a comparison with electronic lab-notebook and
provenance systems.

\section{Conclusion}\label{sec:conclusion}

A small typed language over ordinary set theory, a store machine with a complete ledger, and a handful of Lean theorems about them are enough to
make three frequent errors of long, model-assisted research ill-formed, and to make every step of a record traceable.  They are not enough to
make the record true.  The case study suggests that the combination with pre-registration, append-only correction and independent review is what
makes such a record trustworthy, and that the typing earns its place in that combination as a guard against one specific, common and otherwise
hard to see kind of error.
